\documentclass{article}
\usepackage[T1]{fontenc}
\usepackage{lmodern}
\usepackage{graphicx}
\usepackage{amsmath,amssymb}
\usepackage[a4paper,margin=2cm]{geometry}
\usepackage[absolute,overlay]{textpos}
\setlength{\TPHorizModule}{1mm}
\setlength{\TPVertModule}{1mm}
\usepackage[indonesian]{babel}
\usepackage{hyperref}
\setlength{\emergencystretch}{2em}
% unit: Halaman 55

\begin{document}

% === Halaman 55 (llm) ===

\textbf{Operasi aritmetika pada $\text{GF}(2^m)$}

Misalkan $a = (a_{m-1} \dots a_1 \ a_0)$ dan $b = (b_{m-1} \dots b_1 \ b_0) \in \text{GF}(2^m)$

\begin{itemize}
    \item \textbf{Penjumlahan:} \\
    $a + b = c = (c_{m-1} \dots c_1 \ c_0)$ dimana $c_i = (a_i + b_i) \pmod 2, \ c \in \text{GF}(2^m)$
    
    \item \textbf{Perkalian:} $a \cdot b = c = (c_{m-1} \dots c_1 \ c_0)$ dimana $c$ adalah sisa pembagian polinom $a(x) \cdot b(x)$ dengan \textit{irreducible polynomial} derajat $m, \ c \in \text{GF}(2^m)$
    
    \item Algoritma Rijndael (AES) beroperasi di dalam $\text{GF}(2^8)$
\end{itemize}

\end{document}
